\documentclass[conference]{IEEEtran}
\usepackage{newtxtext,newtxmath}
\usepackage{booktabs}
\usepackage{array}
\usepackage{amsmath}
\usepackage{graphicx}
\usepackage{flushend}
\usepackage[hidelinks]{hyperref}
\usepackage{xcolor}

\newcommand{\modeltag}{\textsuperscript{[model]}}
\newcommand{\vcstag}{\textsuperscript{[VCS]}}
\newcommand{\dctag}{\textsuperscript{[DC/PT]}}
\newcolumntype{L}[1]{>{\raggedright\arraybackslash}p{#1}}

\title{28-nm Capacity-Constrained Typed-Source Execution Islands\\
for Event-Driven Spiking Optical Flow}
\author{Anonymous submission}

\begin{document}
\maketitle

\begin{abstract}
Low firing activity does not automatically accelerate event-driven optical
flow: finite storage, single-port state conflicts, and repeated weight delivery
can dominate useful arithmetic.  We present typed-source execution islands for
a frozen binary-event optical-flow network.  C1 combines exact repeated-product
capture with finite single-port parent state and dead-write suppression.  On
51.84 million source rows, its same-ledger cycle model reduces component time by
40.99\% (1.6945$\times$); a mapped 28-nm nine-SRAM component meets 3-ns
setup/hold and passes 16,549 Formality compare points.  C2 shares Acc24 state across
eight typed-signed sources.  Across five directed workloads, against
equal-bandwidth K1$\times$8, K8 takes 1,913
versus 1,945 directed VCS cycles (1.0167$\times$) and yields 4.5507$\times$
directed-throughput/logic-area with 77.66\% less synthesized logic.  Token-set
broadcast then reuses weight delivery without reusing signed products.  On four
preselected ep34 FC1 G48 service groups, one per DSEC sequence, same-port/cache
RTL reduces post-load VCS cycles from 86,713 to 30,775 (2.8176$\times$) and
execution time by 64.51\%, while weight requests fall by 69.92\%; a matched
logic-only DC ablation adds 0.0118\% area.
Matched hold and power remain open.  C3 closes exact fixed-$T{=}10$ temporal
service at 3 ns.  An operator-local INT8 bridge matches 1,280 integer-oracle
probes, stays within Acc24, and passes a predefined AEE compatibility gate on an
825-frame local DSEC validation set.  We keep model, RTL, and physical evidence
separate and make no whole-network speedup claim.
\end{abstract}

\section{Introduction}
Event cameras make inactivity visible, but sparse arithmetic is only one term
in execution time.  In our optical-flow deployment, nonzero sources still
trigger parent-state traffic, accumulator ownership, and weight-row fetches.
Consequently, a design that reports only operation sparsity can be slower than
an equal-resource zero-skipping baseline.  Prosperity~\cite{prosperity} and
Phi~\cite{phi} demonstrate that useful SNN speedups come from executable reuse
or patterns, not firing rate alone; FireFly-T~\cite{fireflyt} shows the value of
multi-spike service.  Our problem differs in its finite 240-KiB-class working
set, polarity-bearing typed sources, and mixed fixed-$T$ neuron phases.

We make two architectural contributions.  (1) C1 captures exact repeated products subject
to a finite-capacity, single-1RW parent-store contract.  (2) C2 shares
Acc24 and endpoint state across eight typed sources and embeds token-set
broadcast (TSBG), which reuses a weight row while preserving independent signed
products and accumulator contexts.  C3 provides exact fixed-$T{=}10$ temporal
service as network-coverage evidence rather than a third performance claim.
The checkpoint is an evaluation workload rather than a parallel algorithmic
novelty claim.

\section{Workload and Execution Contract}
The frozen Motion C12 ep34 checkpoint obtains 1.199514 AEE on an 825-frame local DSEC validation set and
has a 5.6709\% global firing rate.  Of 105 configured ATLIF wrappers, 12
\texttt{sn2\_q} wrappers are runtime-bypassed and never called, so the capture
invokes 93 (48 $T{=}2$, 45 $T{=}10$).  A separate graph audit finds that 12
invoked \texttt{attn\_sn} return values have no consumer under fixed normal
inference, leaving 81 graph-live services under that fixed call graph
(36 $T{=}2$, 45 $T{=}10$).
All 93 captured ATLIF outputs are binary.  Signed source codes are therefore a downstream
execution-protocol property---for example polarity and correction traffic---not
evidence that the captured ATLIF tensors are analog.
The frozen software path uses Motion-XOR scoring with $\alpha{=}0.125$ and
K as its value carrier, but its configuration has hardware quantization
disabled.  Only the separately evaluated deployment candidate enables Q7 round-to-nearest
scores, a next-power-of-two integer Shiftmax row sum, and Q1.7 gates.
Attention is not promoted to a performance contribution, nor are its
fractional gate values called powers of two.

Each result is assigned one evidence class.  A cycle-model result replays the
same trace and charged resources in software.  A VCS result proves directed
RTL function/protocol properties.  DC/PT/Formality results anchor area, timing,
and equivalence.  Only a result that closes all three classes on one workload
may be called an RTL speedup; the present paper instead reports the three
classes side by side.

The mapping is driven by three exact workload invariants rather than a claim
that every layer has the same sparsity.  Bottleneck convolution exposes
repeated products, but only while parent and destination state remain resident;
this defines C1.  FC tokens can share a weight-row identity while retaining
different signs, destinations, and accumulators; C2 therefore reuses delivery,
not products, through TSBG.  Fixed-$T{=}10$ neurons require a deterministic
state and commit order; C3 preserves that order.  Attention and decoder traces
remain part of the frozen workload identity but are not promoted to additional
performance contributions.

\begin{figure}[t]
\centering
\setlength{\fboxsep}{3pt}
\fbox{\parbox{0.94\columnwidth}{\centering\scriptsize
\textbf{Frozen binary-event workload} $\rightarrow$
\textbf{typed source} $\{\mathrm{row/token},\mathrm{sign},\mathrm{dst},\mathrm{last}\}$\\[3pt]
\begin{tabular}{@{}>{\centering\arraybackslash}p{0.44\columnwidth}|>{\centering\arraybackslash}p{0.44\columnwidth}@{}}
\textbf{C1: Conv product capture} &
\textbf{C2: FC typed-source service}\\
64-row mask directory & four private token contexts\\
exact-subset parent $+$ residual & one weight-row fetch/broadcast\\
dead-write/live-parent test & up to eight signed sources\\
9$\times$128$\times$128-bit 1RW scratch & shared Acc24 fabric\\
\end{tabular}\\[3pt]
\textbf{C3:} exact fixed-$T{=}10$ state/commit}}
\caption{The component islands share a typed-source contract; the branches do
not imply a measured monolithic C1--C2 integration.  TSBG is a C2
weight-delivery specialization.}
\label{fig:overview}
\end{figure}

\section{Microarchitecture}
\subsection{C1: Capacity-Constrained Product Capture}
C1 first builds a ping-pong directory for each 64-row, 16-source mask task.  A
row may select only an earlier exact-subset row as its parent; it then issues
the XOR residual and reconstructs the 96-lane signed product row by adding the
parent.  Stable population-count ordering and a live-parent bitmap make this
relation deterministic and identify final values that no later row can
observe.  Such dead writes are elided exactly.

The execution side combines one earliest-parent lookahead, a two-entry reserved
response queue, and forwarding under a single-1RW deadline-aware arbiter.
Final psum write and row completion are atomic, so backpressure cannot expose a
partially reconstructed row.  The scheduler uses nine 128$\times$128 1RW SRAM
macros and never assumes the ideal simultaneous access that produced the
higher but unphysicalized reuse ceiling.  This is the object difference from
an unbounded product-sparsity opportunity: reuse must survive capacity, port,
parent lifetime, and completion constraints.

For parent mask $\mathcal{M}_p\subset\mathcal{M}_r$, exactness follows from
the disjoint residual rather than value prediction:
\begin{equation}
\mathbf{P}_r=\mathbf{P}_p+
\sum_{s\in\mathcal{M}_r\setminus\mathcal{M}_p}\mathbf{p}_s.
\label{eq:c1exact}
\end{equation}

Table~\ref{tab:c1ladder} exposes the baseline ladder instead of attributing all
inactivity to product capture.  Plain bit skipping barely changes time on this
ledger; exact parent reuse creates the material reduction.  The remaining
12.25\% gap from C1 to the concurrent-access ceiling is the price of the
physical single-port contract, not an omitted claimed speedup.

\begin{table}[t]
\caption{C1 same-ledger capture ladder (51.84 M source rows).}
\label{tab:c1ladder}
\centering
\scriptsize
\begin{tabular}{@{}lrrl@{}}
\toprule
Execution point & Cycles (M) & vs. zero & Boundary \\
\midrule
Strong zero & 648.741 & 1.000$\times$ & baseline \\
Same-coordinate bit & 646.619 & 1.003$\times$ & baseline \\
C1 finite 1RW & 382.849 & 1.6945$\times$ & cycle model \\
Concurrent-access & 341.058 & 1.902$\times$ & ceiling only \\
\bottomrule
\end{tabular}
\end{table}

\subsection{C2: Typed K8 and Token-Set Broadcast}
The K8 fabric accepts at most eight nonzero typed sources and shares the Acc24
array, endpoint, and control rather than replicating eight K1 engines.  A fair
baseline is therefore K1$\times$8 at equal source bandwidth, not one K1 lane.
TSBG changes only weight delivery.  Four token contexts expose a common
source-group row; one fetched row is broadcast to four independent signed
products and Acc24 contexts.  It never reuses a product because equal source
indices can carry different signs or destinations.  Both ordinary and TSBG
modes use the same 192-bit row-live representation and hierarchical priority
encoder, avoiding division/remainder hardware that would bias the physical
comparison.  Bundled event delivery and mini-batch Gustavson execution are
established in ELSA~\cite{elsa}, while SpikeX~\cite{spikex} reuses weights
across time and space.  Our narrower contribution is their mapping to a
polarity-bearing source stream with four private destination/Acc24 contexts
under the same finite row cache and public-port contract.
For a common weight row $W_g$, TSBG shares only its delivery:
\begin{equation}
\operatorname{fetch}(W_g)=1,\quad
A_c\leftarrow A_c+v_cW_g,\quad (v_c,d_c,A_c)\ \text{private}.
\label{eq:tsbg}
\end{equation}

\subsection{C3: Exact Fixed-$T$ Coverage}
C3 schedules the deployed $T{=}10$ ATLIF recurrence and preserves its integer
state and commit order.  It is included to close the temporal service boundary,
not to claim a standalone speedup.

All performance mechanisms have an exact fallback.  If C1 finds no resident
legal parent, it selects the empty parent and issues the original source row.
If token contexts disagree on a weight-row identity, C2 emits separate fetches;
sign, destination, and Acc24 state are never coalesced.  Backpressure changes
only service order, because the typed-source terminal bit and atomic commit keep
ownership explicit.  Thus sparsity affects delivered work without changing the
operator result.

\section{Evaluation}
All physical results use a commercial 28-nm flow at 3.0 ns.  The reported
points are prelayout, use ideal clocks and ZeroWireload, and have no extracted
parasitics; C1 area includes nine SRAM macros, whereas C2 and C3 are logic-only.
The component
comparison in Table~\ref{tab:main} deliberately keeps cycle, area, and proof
scope visible.  C1 model cycles replay ten ep34 samples and 51.84 million source
rows; per-sample ratios span 1.661--1.723$\times$, with a 1.6946$\times$
geometric mean.  C2 VCS cycles use five frozen directed workloads.  TSBG's
four-sequence VCS calibration uses the first captured sample of each sequence,
layer 28, and tokens 0--3; the selection is fixed independently of performance.
The broader software screening uses 960 FC1/FC2 pairs from 40 samples and four
DSEC sequences.  Task accuracy binds the frozen checkpoint to a deployment subset
on the same population and metric aggregation.  Because the historical baseline
and candidate use different GPU backend flags, their difference is
baseline-referenced rather than a single-factor causal estimate.  We do not
claim full-network INT8 or a whole-network RTL speedup.

\begin{table*}[t]
\caption{Admitted component evidence; ``open'' claims are not inferred.}
\label{tab:main}
\centering
\scriptsize
\begin{tabular}{@{}L{0.055\textwidth}L{0.235\textwidth}L{0.275\textwidth}L{0.365\textwidth}@{}}
\toprule
Line & Performance & Physical anchor & Evidence boundary \\
\midrule
C1 & 1.6945$\times$; $-40.99\%$ time\modeltag & 166,514 $\mu$m$^2$; PT setup/hold +27.9/+1.8 ps; 16,549 compare points pass & Model ratio plus a separate mapped physical anchor; one real-mask tile is event-count calibrated in VCS \\
C2 & 1,913 vs 1,945 directed cycles (1.0167$\times$)\vcstag; 4.5507$\times$ directed-VCS throughput/DC logic-cell area & logic-only 130,823 vs 585,535 $\mu$m$^2$; $-77.66\%$ & K8 vs equal-bandwidth K1$\times$8; matched hold/power open \\
C2/TSBG & 86,713 vs 30,775 G48 execute cycles (2.8176$\times$)\vcstag; time $-64.51\%$; weight requests $-69.92\%$ & matched logic-only 249,710 vs 249,740 $\mu$m$^2$; $+0.0118\%$; setup met & Four fixed ep34 real-activity groups; identical 383-cycle preloads excluded; no full-FC/system claim \\
C3 & 17 cycles/tile & 63,756 $\mu$m$^2$\dctag; DC/PT setup+hold met; 11,180 compare points pass & Exact service; Formality compares pre- and post-hold-repair mapped netlists, not RTL-to-gate; no speedup/energy claim \\
\bottomrule
\end{tabular}
\end{table*}

\begin{table}[t]
\caption{Frozen-population local DSEC accuracy against a historical,
backend-mismatched baseline (825 frames; lower is better).}
\label{tab:accuracy}
\centering
\scriptsize
\begin{tabular}{@{}lrrr@{}}
\toprule
Metric & Baseline & Candidate & $\Delta$ \\
\midrule
AEE & 1.199514 & 1.197367 & $-0.002147$ \\
AAE & 5.400641 & 5.412808 & $+0.012167$ \\
AAE-Bench. & 5.106363 & 5.121619 & $+0.015256$ \\
DSEC-Fl & 5.313360 & 5.328834 & $+0.015475$ \\
\bottomrule
\end{tabular}
\end{table}

AEE (average endpoint error) is the preselected compatibility-gate metric.
AAE is the legacy 2-D angular error; AAE-Bench. is the Middlebury/Barron 3-D
angular error; DSEC-Fl is the percentage with endpoint error above 3 pixels and
5\% of ground-truth flow magnitude.

Table~\ref{tab:accuracy} uses the same 825 frames, 18 sequences, and
48,152,523 valid pixels.  The candidate applies hardware-order attention to all
12 blocks and dyadic-INT8 QDQ to four C1 Conv3x3 and four decoder
ConvTranspose weights; every transformed weight is reached once per frame.
Other operators remain at checkpoint precision.  The historical baseline
enabled TF32/cuDNN benchmarking whereas the candidate disables both; therefore
the table demonstrates accuracy-gate compatibility, not a causal accuracy
improvement from quantization.
Aggregate AEE falls by 0.002147, but the three auxiliary errors in
Table~\ref{tab:accuracy} rise and 10 of 18 sequences regress in AEE; we claim
only frozen-population gate compatibility, not universal improvement.

The C1 mapped component includes its nine SRAM leaves.  It dissipates 29.08 mW
and 22.07 nJ in a separately bounded 64-row, 253-cycle directed window.  This is
a mixed-corner, prelayout estimate without SPEF, not energy per frame.  C2 K8
and K1$\times$8 meet setup in fresh matched synthesis; independent matched hold
and power remain outside the admitted row.

For C1, a separately sealed ep34 tile uses the first 64 real support masks of
the cycle ledger.  VCS matches 196 accepted issues, 58 parent edges, 31
dead-write elisions, 54/33 macro reads/writes, four forwards, six deadline
holds, 14 stalls, and 64 commits/completions.  Its lane values are deterministic
synthetic signed-12-bit data with zero prior psum; this calibrates modeled event
counts and directed arithmetic, not the 1.6945$\times$ cycle ratio itself.

For numerical binding, all eight retained Conv3x3/ConvTranspose weights are
quantized per output channel to narrow-range signed INT8 with a power-of-two
scale.  An independently checked export has zero pre-clip violations; the
largest absolute-code accumulator bounds are 200,219 for C1 and 87,136 for the
decoder, both below Acc24.  The subsequent operator-local replay covers 160
full calls and 3,271,680,000 output elements.  Against FP32 it gives 0.009106
MAE, 0.012407 RMSE, 0.239204 maximum absolute error, and 0.9999789 cosine
similarity.  All 1,280 independently sampled Python-integer probes match;
the full observed final-accumulator range is $[-29{,}680,27{,}619]$, and the
formal worst-case prefix bound is 200,219, both within signed Acc24.  These are
operator-local numerical results, not task AEE, downstream-neuron equivalence,
or a hardware speedup.  Table~\ref{tab:accuracy} supplies the separate task-level
accuracy binding without promoting this operator-local bridge to a full-network
integer implementation.
As a C2 scheduling study, TSBG's full-capture software replay reduces serialized
modeled cycles from 29.41 billion to 11.61 billion and weight traffic by
65.20\%\modeltag, with a minimum per-sequence ratio of 2.517$\times$.  A
separate VCS calibration fixes, before measuring performance, the first sample
from each sequence, layer 28's first 48-source-group FC1 service, and tokens
0--3.  Across these four groups, ordinary LRU4 takes 86,713 post-load execute
cycles versus 30,775 for TSBG-B4 (2.8176$\times$); bundle/scalar weight requests
fall from 7,620/60,960 to 2,292/18,336 ($-69.92\%$).  The fixture supplies real
ep34 activity masks; all nonzero codes in the measured groups are $+1$.
Deterministic INT8 weights exercise arithmetic but do not affect scheduling.
The same runs check bank stalls/reordering and stale responses; a subsequent
synthetic recovery phase checks signed products, reset, and the $-128$ corner.  This is
an RTL component microbenchmark, not full-FC or system speedup.  The physical
comparison elaborates the same B4 RTL with only
\texttt{SCHEDULE\_MODE=0/1} changed; row-live representation, cache capacity,
public ports, and synthesis constraints are otherwise identical.
In that matched logic-only ablation, ordinary LRU4 and TSBG-B4 occupy
249,710 and 249,740 $\mu$m$^2$, respectively, a 0.0118\% increment, and both
meet 3-ns setup.  Both have a diagnostic hold slack of $-16.4$ ps.  State is
implemented with standard cells in this experiment; power, SRAM-macro cost,
and hold closure remain open.

Table~\ref{tab:tsbgseq} shows both the full-capture model and the fixed G48 RTL
calibration; neither is promoted to full-network speedup.

\begin{table}[t]
\caption{TSBG-B4 full-capture model and fixed G48 RTL calibration.}
\label{tab:tsbgseq}
\centering
\scriptsize
\begin{tabular}{@{}lrr@{}}
\toprule
DSEC trace & Full model & Fixed G48 VCS \\
\midrule
interlaken\_01\_a & 2.544$\times$ & 2.858$\times$ \\
thun\_01\_b & 2.550$\times$ & 2.979$\times$ \\
zurich\_city\_09\_a & 2.525$\times$ & 2.681$\times$ \\
zurich\_city\_12\_a & 2.517$\times$ & 2.743$\times$ \\
\textbf{aggregate} & \textbf{2.534$\times$} & \textbf{2.818$\times$} \\
\bottomrule
\end{tabular}
\end{table}

\section{Related Work and Comparison Discipline}
Prosperity reports average accelerator speedup against PTB and separately
isolates product over bit sparsity~\cite{prosperity}; the latter is the relevant
mechanism-level comparison to C1.  Phi chooses its storage point by DSE and
compares its complete hierarchical-pattern architecture with published
baselines~\cite{phi}.  We adopt the useful evaluation pattern---cycle simulation
plus RTL/physical anchors and a baseline ladder---without relabeling their
numbers as ours.  ELSA and SpikeX establish bundle/Gustavson and cross-window
weight-reuse priors~\cite{elsa,spikex}; TSBG is therefore claimed only as a
typed-signed, finite-context delivery specialization inside C2.  FireFly-T's
multi-spike decoder motivates an equal-service comparison for
C2~\cite{fireflyt}.  Lossy pruning is outside the scope of this exact
component paper.

SpinalFlow sorts compressed spike streams to preserve useful state locality,
while SNE makes convolutional work proportional to the incoming event
stream~\cite{spinalflow,sne}.  They establish event-driven execution, but do not
remove C1's repeated products under a finite 1RW parent lifetime or C2's repeated
weight delivery across private signed contexts.  A recent 28-nm event-based
optical-flow chip jointly reports operations, external-memory traffic, latency,
energy, and AEE~\cite{ciccflow}.  We follow that separation of metrics, but do
not numerically rank its silicon measurements against our prelayout component
anchors.

For a co-designed network, a published chip rarely supports the identical graph.
We therefore compare at three levels: same-network component baselines,
iso-workload public-artifact opportunities, and published technology-normalized
metrics.  Results from different levels are never multiplied.

\begin{table}[t]
\caption{Mechanism boundary; prior numbers are not relabeled as ours.}
\label{tab:prior}
\centering
\scriptsize
\begin{tabular}{@{}L{0.20\columnwidth}L{0.21\columnwidth}L{0.49\columnwidth}@{}}
\toprule
Work & Prior object & Boundary retained here \\
\midrule
Prosperity~\cite{prosperity} & repeated products & C1 adds finite single-1RW parent lifetime, exact residual reconstruction, and atomic completion \\
Phi~\cite{phi} & hierarchical patterns & no substitution for C1 dynamic parents or C2 signed contexts \\
FireFly-T~\cite{fireflyt} & multi-spike service & motivates equal-bandwidth K8, not a one-lane baseline \\
ELSA~\cite{elsa}/\allowbreak SpikeX~\cite{spikex} & bundle/weight reuse & TSBG preserves private sign, destination, and Acc24 state \\
SpinalFlow~\cite{spinalflow}/\allowbreak SNE~\cite{sne} & event-stream execution & do not subsume finite-parent product reuse or typed cross-token delivery \\
CICC optical flow~\cite{ciccflow} & silicon system metrics & metric template only; no cross-platform speedup against prelayout islands \\
This work & exact parent + typed delivery & same-ledger model, VCS, and 28-nm islands; no system-speedup claim \\
\bottomrule
\end{tabular}
\end{table}

\subsection{Limitations}
The C1 ratio is a cycle-model result calibrated by one real-mask RTL tile; its
mapped nine-SRAM component is a separate physical anchor.  C2's area-efficiency
comparison is logic-only.  TSBG's RTL cycle result covers four fixed G48 groups,
while its full-capture external weight-service time remains modeled.  We leave
matched-macro C2 hold/power, full-network FPS/energy, and monolithic island
integration open instead of deriving them from component ratios.

\section{Conclusion}
The principal hardware result is not firing sparsity itself, but its executable
capture under finite storage and service constraints.  C1 provides bounded
exact product reuse; C2 turns typed multi-source service into area efficiency
and uses TSBG to suppress repeated weight delivery; C3 closes exact temporal
service.  Separating model, VCS, and physical evidence yields a compact ISCAS
component paper without an unsupported whole-network claim.

\bibliographystyle{IEEEtran}
\bibliography{references}
\end{document}
